\clearpage
\begingroup
\fancyhead[L]{\small THEORIE DER BEWEGUNG DES MONDES.}
\renewcommand{\headrulewidth}{0pt}
\fancypagestyle{VIIPrinted634}{
  \fancyhead{}
  \fancyhead[L]{\small 634}
  \fancyhead[C]{\small NACHLASS.}
}
\SourceLeaf{640}\thispagestyle{VIIPrinted634}
\begin{center}
\textbf{FÜNFTER ABSCHNITT.}\\[8pt]
\textit{Zweite Annäherung zur Länge.}
\end{center}
\par\noindent\hspace*{\parindent}Wir nehmen zuvörderst die 2\textsuperscript{te} Fundamentalgleichung vor und entwickeln
\[
\begin{aligned}
\frac{\mathrm du^2}{\mathrm dv^2}\pi
={}&\tfrac32nn\frac{\mathrm du^2}{\mathrm dv^2}\frac{B^3}{R^3}
\Bigl\{\sin2(v-V)\\
&+\frac{1}{Rp}\bigl(\tfrac14\sin(v-V)+\tfrac54\sin3(v-V)\bigr)\Bigr\}.
\end{aligned}
\]
\noindent Setzen wir hier \(\dfrac{1}{Rp}\) als eine Constante an und bezeichnen dieselbe durch \(7\),
so lässt sich der zweite Factor dieses Ausdrucks in [die Form bringen:]
\[
\begin{aligned}
&(1-4\varepsilon\varepsilon)\sin2E+\tfrac14\,7\sin E+\tfrac54\,7\sin3E\\
&\quad+2\varepsilon\sin(2E+a)-2\varepsilon\sin(2E-a)\\
&\quad+\tfrac34\varepsilon\varepsilon\sin(2E+2a)+\tfrac{13}{4}\varepsilon\varepsilon\sin(2E-2a)\\
&\quad-2ne\sin(2E+M)+2ne\sin(2E-M)\\
&\quad+\tfrac34nee\sin(2E-2M)+\tfrac14ngg\sin(2E-2\lambda)\\
&\quad\bigl[-\tfrac34nee\sin(2E+2M)-\tfrac14ngg\sin(2E+2\lambda)\bigr];
\end{aligned}
\]
\noindent folglich
\[
\left[\frac{\mathrm du^2}{\mathrm dv^2}\pi\right]
=\tfrac32nn(1+2ee+\tfrac32\varepsilon\varepsilon)\times
\]
\[
\begin{array}{r|l}
\multicolumn{1}{c}{\sin}&\\
2E&1-4\varepsilon\varepsilon\\
E&\tfrac14\,7\\
3E&\tfrac54\,7\\
4E&-\tfrac{11}{4}nn\\
2E+M&2e-2ne\\
2E-M&2e+2ne\\
2E+2M&\tfrac52ee-\tfrac{19}{4}nee
\end{array}
\]
\[
\begin{array}{r|l}
\multicolumn{1}{c}{\sin}&\\
2E-2M&\tfrac52ee+\tfrac{19}{4}nee\\
4E-M&\tfrac{15}{4}ne\\
2E+a&\tfrac12\varepsilon\\
2E-a&-\tfrac72\varepsilon\\
2E+2a&0\\
2E-2a&\tfrac{17}{2}\varepsilon\varepsilon\\
2E+M+a&e\varepsilon
\end{array}
\]
\fancypagestyle{VIIPrinted635}{
  \fancyhead[L]{}
  \fancyhead[C]{\small THEORIE DER BEWEGUNG DES MONDES.}
  \fancyhead[R]{\small 635}
}
\clearpage\SourceLeaf{641}\thispagestyle{VIIPrinted635}
\[
\begin{array}{r|l}
\multicolumn{1}{c}{\sin}&\\
2E+M-a&-7e\varepsilon\\
2E-M+a&e\varepsilon\\
2E-M-a&-7e\varepsilon\\
M&\tfrac{15}{4}ne\\
2E+2\lambda&\tfrac12gg-\tfrac14ngg\\
2E-2\lambda&\tfrac12gg+\tfrac14ngg
\end{array}
\]
\[
\begin{array}{r|l}
\multicolumn{1}{c}{\sin}&\\
2E+a-2\lambda&\tfrac14gg\varepsilon\\
2E-a-2\lambda&-\tfrac74gg\varepsilon\\
a&0\\
\multicolumn{2}{c}{\left[\begin{array}{r|l}
2E+2M+a&\tfrac54ee\varepsilon\\
2E-2M-a&-\tfrac{35}{4}ee\varepsilon
\end{array}\right]}
\end{array}
\]
\noindent Folglich gibt die Integration
\[
\log pp\frac{\mathrm du}{\mathrm dv}=\log A+
\]
\[
\begin{array}{r|l}
\multicolumn{1}{c}{\cos}&\\
2E&-\tfrac34nn-\tfrac34n^3-\tfrac34n^4-\tfrac32nnee+\tfrac{15}{8}nn\varepsilon\varepsilon\\
4E&\tfrac{33}{32}n^4\\
E&-\tfrac38nn\,7-\tfrac38n^3\,7\\
3E&-\tfrac58nn\,7-\tfrac58n^3\,7\\
2E+M&-nne+\tfrac13n^3e\\
2E-M&-3nne-9n^3e\\
2E+2M&-\tfrac{15}{16}nnee\\
2E-2M&\tfrac{15}{8}nee+\tfrac{159}{32}nnee\\
4E-M&-\tfrac{15}{8}n^3e\\
2E+a&-\tfrac38nn\varepsilon-\tfrac{3}{16}n^3\varepsilon\\
2E-a&\tfrac{21}{8}nn\varepsilon+\tfrac{63}{16}n^3\varepsilon\\
2E+2a&0\\
2E-2a&-\tfrac{51}{8}nn\varepsilon\varepsilon
\end{array}
\]
\[
\begin{array}{r|l}
\multicolumn{1}{c}{\cos}&\\
2E+M+a&-\tfrac12nne\varepsilon\\
2E+M-a&\tfrac72nne\varepsilon\\
2E-M+a&-\tfrac32nne\varepsilon\\
2E-M-a&\tfrac{21}{2}nne\varepsilon\\
M&-\tfrac{45}{8}n^3e\\
2E+2\lambda&-\tfrac{3}{16}nngg\\
2E-2\lambda&\tfrac38ngg-\tfrac{3}{32}nngg\\
2E-2\lambda+a&\tfrac38ngg\varepsilon\\
2E-2\lambda-a&-\tfrac78ngg\varepsilon\\
a&0\\
\multicolumn{2}{c}{\left[\begin{array}{r|l}
2E-2M+a&\tfrac{15}{8}nee\varepsilon\\
2E-2M-a&-\tfrac{35}{8}nee\varepsilon
\end{array}\right]}
\end{array}
\]
\par\noindent\hspace*{\parindent}Hieraus folgt sogleich der Werth von \(p^4\dfrac{\mathrm du^2}{\mathrm dv^2}\), welchen wir in der ersten
Gleichung zu substituiren haben; er ist nemlich
\[
=A(1+\text{Aequatt. praecc. dupl.}+\tfrac{9}{16}n^4+\tfrac{9}{16}n^4\cos4E).
\]
\noindent Ferner findet sich leicht aus der oben gefundenen Formel für \(\theta\):\par
\vfill
\noindent\hfill80\textsuperscript{*}\par
\clearpage
\endgroup
